\section*{摘要}
\addcontentsline{toc}{section}{摘要}
本实验以一台 TonyPi 人形机器人和四台 TurtleBot2 为平台，研究\textbf{抱物穿行与多机器人协同跟随}。任务要求 TonyPi 抱持目标物体完成左、右转向并在终点放置，移动机器人保持跟随并满足避碰、目标保持和终点停靠要求。本报告重点阐释第一小组承担的\textbf{TonyPi 路线控制}与\textbf{方式二建图避障跟随}。

TonyPi 采用\textbf{动作组执行与视觉反馈相结合}的分层结构，以\textbf{抓取末帧锁定上肢姿态}并复用下肢步态。最新路线通过\textbf{A/B/C 自定义固定标记识别与 PnP 位姿解算}取得距离和方位，利用\textbf{分层状态机}组织抱起、走向 A、左转寻找 B、走向 B、右转寻找 C、走向 C及放置，并采用\textbf{多帧中位数观测、自适应纠偏与有界异常处置}控制离散动作的切换。

方式二以\textbf{GMapping 建图与 AMCL 定位}建立场地参考系，将\textbf{AR 目标观测按采样时刻变换至地图坐标}，沿\textbf{目标历史轨迹的弧长回退}生成跟随点，再通过\textbf{move\_base/DWA 路径规划与局部避障}输出运动指令。\textbf{上层有效性检查与独立速度门控}共同约束任务状态、目标、定位、TF、地图和扫描时效，实现异常情况下的取消目标与零速输出。

报告以源码、动作数据库和已有运行记录为依据，区分算法实现、离线验证、局部实机运行和完整任务验收。当前证据尚不足以证明全路线及全队实验完成，文中明确说明\textbf{最新人形入口与任务心跳的集成边界}，不将默认参数、解析曲线或合成输入作为实测成绩。

\noindent\textbf{关键词：}TonyPi；有限状态机；视觉标记；地图导航；轨迹跟随；多机协同
